\section{Introducción}

El precio del combustible tiene efectos sobre el bienestar social tanto por la asignación eficiente de los recursos mediante precios competitivos como por su incidencia en los costos de la cadena productiva. En Chile, los mecanismos institucionales determinan los costos de distribución minorista, y por ende los márgenes de venta son el componente principal por el cual se expresa la competencia entre estaciones de servicio (EDS). Medir y caracterizar la respuesta frente a un nuevo competidor orienta la política de competencia respecto al rol de las barreras de entrada.

El propósito de este estudio es investigar el efecto disciplinador inducido por un nuevo competidor en un mercado particularmente concentrado y con altas barreras de entrada. El informe de la \textcite{FNE2025informe} caracterizó el mercado de distribución minorista de combustible en Chile bajo una alta concentración e integración vertical en tres franquicias que controlan alrededor del $80\%$ de las EDS. El informe detalla que la regulación dentro de toda la cadena productiva elevan las barreras de entrada dificultando la presencia de nuevos competidores. Adicionalmente, el mercado es naturalmente de caracter geográfico, los márgenes de venta que una EDS de servicio puede imponer dependen no solo del número de competidores sino de la distancia entre ellos.

Considerando estas características, el mercado de combustible resulta interesante de investigar puesto que la alta visibilidad del precio y la perfecta homogeneidad del producto aislan el componente espacial como uno de los más importantes en la determinación del precio. Dado esto, el mercado de combustibles ha sido investigado bajo esta dimensión en múltiples trabajos, principalmente de corte transversal midiendo la densidad de competidores \parencite{CLEMENZGUGLER2006}. De manera estructural, \textcite{HOUDE2012} especifica los trayectos de los consumidores, y no su lugar de residencia, como las ubicaciones de un modelo de tipo Hotelling, de modo que la diferenciación espacial pasa a depender de la red vial y de la dirección de los flujos de tráfico. En Chile existen dos precedentes cercanos, ambos tesis de magíster de la Universidad de Chile: \textcite{CASTILLO2018} mide el efecto de la cercanía de estaciones rivales sobre el precio, distinguiendo entre redes integradas e independientes y construyendo distancias en tiempo de viaje entre todas las EDS georreferenciadas del país, y \textcite{SOLIS2025} documenta la dispersión de precios entre estaciones cercanas y la atribuye a costos de búsqueda del consumidor. En esta oportunidad se busca capturar el efecto promedio de una entrada sobre el precio de las incumbentes en un mercado local particular, siguiendo los métodos de \textcite{FISCHER2025, KOH2022}, que serán los articulos seminales de esta investgación. 

Una diferencia importante con respecto a estos articulos que justifica la exploración de este tema en Chile es que los mecanismos institucionales tales como el Mecanismo de Estabilización de Precios de los Combustibles (MEPCO) son públicos y de frecuencia semanal, al igual que los precios de paridad publicados por la Empresa Nacional del Petróleo (ENAP), lo que permite aproximar el costo mayorista que enfrentan las estaciones con información pública disponible.\footnote{\textcite{KOH2022} no observa márgenes y por tanto restringe su análisis a precios de venta y aproximaciones de márgenes que tienen una relevancia acotada. \textcite{FISCHER2025} compra la información de márgenes comercialmente. El MEPCO por lo tanto es útil en poder aproximar estos márgenes de venta con información pública} 

Como insumo para la investigación se explotan los datos históricos de las actualizaciones de precio de cada EDS junto a su localización y distribuidor. Adicionalmente, para aproximar el margen de venta de las estaciones se utilizan datos del precio mayorista de las estaciones considerando el MEPCO. Con el fin de falsear las hipótesis que se exponen más adelante se utiliza un estudio de eventos mediante un estimador de diferencias en diferencias con efectos fijos. El método encuentra efectos causales bajo el supuesto de exogeneidad del timing de entrada desde el punto de vista de la incumbente \parencite{ARCIDIACONO2020}, el cual se defiende más adelante.

Esta investigación aporta a la literatura en tres aspectos. El primero es de medición. \textcite{KOH2022} no observan el precio mayorista y deben restringir su análisis a los precios de venta, \textcite{FISCHER2025} acceden a márgenes a través de datos comerciales. El MEPCO permite aproximar el margen bruto de cada estación con información enteramente pública, y el trabajo documenta cómo hacerlo. 

El segundo aporte es testear mediante la estrategia empírica si los márgenes ex-ante de la entrada determinan el tamaño del efecto displicanador. Se encuentra que las estaciones que tienen márgenes mayores con respecto a su comuna son quienes más reducen sus precios ante la entrada de una competidora. 

El tercer aporte es reportar nueva evidencia causal sobre el comportamiento de las EDS frente a la entrada de nuevos competidores en una economía en desarrollo. Gracias a estas estimaciones tenemos más información para delimitar de manera más precisa los mercados locales en el mercado de combustibles chileno, relevante para un análisis de libre competencia. 